\section{Who may amend, and where entitlement comes from}
\label{sec:14.2}
\runin{Who may amend}

At about a quarter past three in the morning of 14 May 2025, someone at xAI changed the instructions its chatbot Grok ran under on X. For hours afterward Grok answered questions about streaming services and sinus remedies with claims about “white genocide” in South Africa. The company called the change unauthorized, and promised to publish its prompts and review changes to them \autocite{cnbc2025grok}. Nothing in the artifact told the edit apart from a correction: a change to what a system will and won't do is the same operation whether it repairs the floor or removes it.

The floor's real structural position isn't the judge's but the unpublished administrative disposition — unreasoned, unindexed, non-precedential, affected parties unable to find each other, the institution the only repeat player. The comparison with a constitutional court is instructive because binding force is what makes an error \emph{correctable}: it makes the error recur, recurrence produces litigants, and disagreement across jurisdictions forces a higher court to settle it. A floor's error is just as uniform and recurs just as reliably, and nothing converts recurrence into visibility.

So there is no separate amendment procedure to design: the custody terms are the procedure. The procedure has conditions of its own: more parties must agree to a change than to run the deployment; the attempt reaches the record whether it carries or fails (the failed attempt is more informative — a floor nobody tried to change and one whose removal was refused look identical from outside); and an amended floor binds forward only, so a retrained model can't be offered as evidence the act was never covered. The cost of binding forward only is that nobody harmed while the wrong version ran gets anything from the fix. A finding that a refusal was wrongful is made on the same terms: by someone other than the party that wants the change, on a record, in a form others can read. So the appeal route for a wrongful refusal and the amendment procedure are one piece of machinery, and the two uses differ only in who may start it.

It is harder than the constitutional case, because an amendment takes effect in a deployment only if it's \emph{pushed}, so the procedure necessarily retains a channel the operator could use for silent revision. Closing it would mean no deployment could ever be corrected.

\runin{Where entitlement comes from}

The authority of the floor's contents doesn't come from their having been written down; it comes from feedback. The American instance shows why: the Declaration's authors struck the clause condemning the slave trade at two delegations' insistence, and what remained was general enough to be turned on them — as Black petitioners in Massachusetts were already doing, in the Declaration's own terms, within six months \autocite{loc2010declaration,hall1777petition}. The legislature did nothing. A war and three amendments made the proposition binding. The time that took measures how badly the channel was blocked. The record bound no one; it told the enforcement where to aim. The petitioners were the excluded population, and they had what the people a floor protects usually lack: the text in hand, and knowledge of what it left out. The people a floor protects mostly don't know the deployment exists, and have no notice, no standing and no route. So those positioned to hold the floor's author to what it left out aren't the protected population — they're the people who have to apply it.

That constituency is defined by the act. Anyone fine-tuning open weights, writing a system prompt or configuring a product's refusals has to translate the floor for a system of their own, and meets its omissions first, so the population measuring the floor is already larger than the frontier labs and grows as building gets cheaper. What shrinks with it is leverage: the floor binds a lab partly because the lab is visible, capitalized and reachable by a court, and a thousand laptops are legible to nobody. I rely for legibility on the concentration the institutional chapters treat as the disease.
